% =====================================================================
% Results — efficiency. Phase 3 of RESEARCH_PLAN.md.
% Tables are checked-in, hand-audited artifacts; paper/build.sh compiles them.
% =====================================================================
\section{Results: Efficiency}
\label{sec:results-efficiency}

The retained SDPA measurements report first-token-time (FTT, prefill latency,
CUDA-event timed) and time-per-output-token. They do not include an end-to-end
latency column.
Every cell is $\geq 5$ repeats with warmup discarded; Table~\ref{tab:t3}
shows medians, while repeat-level rows remain in the local source store.
The earlier hypothesis timing that appeared alongside no seed control and
one configuration does not survive into this section. Realized speedup is
computed against the \texttt{full} arm on identical shapes.

\subsection{SDPA timing versus sparsity}
\label{sec:results-efficiency-analytical}

Table~\ref{tab:t3} reports median TTFT, TPOT, and analytical mask sparsity at
$\ell=512$ and $Q=64$ under SDPA. Full/DDI TTFT ratios remain near one even as
sparsity rises, so the additive mask does not realize the analytical pair-count
reduction on this backend.

\begin{table}[t]
\centering
\caption{T3 --- SDPA TTFT, TPOT, and DDI mask sparsity by document count $N$
    at $\ell=512$ and $Q=64$. Values are hand-audited from the local Phase~3
    timing store; no end-to-end latency is shown.}
\label{tab:t3}
\IfFileExists{generated/T3.tex}{%
    \begin{tabular}{rccccccr}
\toprule
 & \multicolumn{3}{c}{TTFT (ms)} & \multicolumn{3}{c}{TPOT (ms)} & \\
$N$ & full & ddi\_qf & blind & full & ddi\_qf & blind & sparsity \\
\midrule
1 & 60.35 & 61.59 & 61.39 & 26.03 & 24.45 & 24.09 & 0.00 \\
3 & 151.03 & 154.69 & 154.10 & 24.43 & 24.26 & 24.06 & 0.61 \\
5 & 250.31 & 253.06 & 253.40 & 26.32 & 24.34 & 24.31 & 0.76 \\
9 & 485.39 & 486.65 & 486.95 & 25.62 & 25.64 & 25.64 & 0.86 \\
13 & 765.89 & 766.27 & 765.43 & 30.18 & 30.21 & 30.21 & 0.91 \\
25 & 1900.95 & 1891.02 & 1890.71 & 43.62 & 43.63 & 43.66 & 0.95 \\
\bottomrule
\end{tabular}
%
}{%
    \todoreport{T3 realization table}{make timing-N}%
}
\end{table}

\subsection{Archived realization figure}
\label{sec:results-efficiency-realization}

The checked-in F3 PDF is not included because its matching raw Flex timing log
is absent and it does not match the output of the retired renderer. Recreating
the figure from preserved evidence is part of VIBEQUEUE-50.

\subsection{On the realization gap}
\label{sec:results-efficiency-gap}

A gap between analytical and realized speedup is expected --- the
analytical count in Section~\ref{sec:method-cost} is a pair count, not a
wall-clock model --- and quantifying its size is a first-class result of
this study, not a failure to explain. Candidates for the gap fall into a
short list. Per-block kernel launch overhead is dominant at small $\ell$
and dilutes with larger blocks. Memory bandwidth caps saturate the
sub-quadratic regime and prevent the compute reduction from translating
proportionally. The query block's $Q \cdot N \ell$ term
(Equation~\ref{eq:cost-qf}) dominates once $N$ is small, so the
analytical speedup is itself modest in that regime. And end-to-end,
prefill may simply not be the bottleneck at the shapes measured, so a
prefill-only win dilutes to near-nothing in real serving.

\paragraph{Measurement to date: structural on \texttt{sdpa}.}
The realization ratio measured on the primary instruct model at
$Q = 64$, $L = 512$ under the \texttt{sdpa} backend collapses from
$\approx 0.98$ at $N = 1$ to $\approx 0.06$ at $N = 25$, while
analytical speedup grows monotonically from $1.0\times$ to
$\approx 16.6\times$ and mask sparsity climbs to $\approx 94\%$.
Realized wall-clock speedup remains pinned near $1.0\times$ throughout.
The picture that emerges is that the analytical/realized \emph{gap
itself grows} with $N$ under this kernel: DDI's block-diagonal sparsity
does not convert to wall-clock speedup at any tested shape with
\texttt{sdpa}.

The root cause is structural, not incidental: passing a 4D additive
attention mask to \texttt{sdpa} forces it off the fused flash/efficient
kernel and onto the slow math backend, which materializes the full
$T \times T$ attention matrix regardless of sparsity. The mask that
should skip computation instead adds overhead. FlashAttention-2 was
scoped as the kernel-fixable alternative but could not be compiled
on the test platform (system \texttt{nvcc} 12.2 cannot compile
torch 2.13 headers; no prebuilt wheel exists for Python 3.13).

\subsection{Realized speedup via \texttt{flex\_attention}}
\label{sec:results-efficiency-flex}

PyTorch 2.13's \texttt{flex\_attention}
\cite{flex_attention} compiles a block-sparse attention kernel from a
mask specification, skipping masked $\mathbf{Q}\mathbf{K}^\top$ blocks
at the kernel level rather than materializing and masking the full
matrix. We integrated it as a third backend: each Qwen2 attention layer
is patched to use the compiled \texttt{flex\_attention} kernel with a
\texttt{BlockMask} built from the DDI document boundaries, while decode
(Q\textsubscript{LEN}=1, memory-bound) falls back to \texttt{sdpa}.

Table~\ref{tab:t3-flex} retains archived batch-1 prefill timing observations
under \texttt{flex\_attention} alongside the \texttt{sdpa} baseline at
matching shapes. They are not an admitted end-to-end throughput result: the
current full-model parity gate fails at the smaller proxy shape
$N{=}2$, $\ell{=}128$, $Q{=}64$. Cosine similarity/token-argmax agreement is
$0.993342/87.19\%$ for \texttt{full}, $0.998553/95.00\%$ for
\texttt{ddi\_qf}, and $0.998701/94.38\%$ for \texttt{blind}, below the
required 0.999 and 100\% gates. In particular, they cannot support a
batched-serving claim. Two implementation observations remain:

\begin{enumerate}
\item \textbf{\texttt{sdpa} timing shows the regression.} DDI speedup
  (full\,/\,ddi\_qf) under \texttt{sdpa} is flat at 0.98--1.01$\times$
  across all shapes --- the 4D mask penalty exactly cancels the sparsity
  benefit.
\item \textbf{Archived \texttt{flex\_attention} timing shows a prefill
  signal, not a validated model result.} The batch-1 DDI timing ratio under
  \texttt{flex} grows with $N$ and
  $\ell$: from 1.05$\times$ at $N{=}9$, $\ell{=}512$ to
  \textbf{1.29$\times$} at $N{=}9$, $\ell{=}2048$ (88\% sparsity).
  The total system speedup --- \texttt{flex+ddi\_qf} versus
  \texttt{sdpa+full} --- reaches \textbf{2.13$\times$} at that shape.
   The original \texttt{create\_block\_mask} implementation exhausts A100
   memory while constructing the $N{=}25$, $\ell{=}2048$ mask representation.
   A direct sparse \texttt{BlockMask} construction removes that allocation,
   but its full-model parity gate fails in the available CUDA/Dynamo runtime.
   That failure occurs before any batched timing can be admitted, so neither
   the target cell nor a smaller-shape batch-throughput claim is reported.
\end{enumerate}

\begin{table}[t]
\centering
\caption{T3-flex --- archived batch-1 prefill timings: \texttt{sdpa} vs
    \texttt{flex\_attention}. DDI ratio = full\,/\,ddi\_qf within each
    backend. System speedup = \texttt{sdpa+full}\,/\,\texttt{flex+ddi\_qf}.
    $Q=64$, Qwen2.5-7B-Instruct, A100 80GB. The largest DDI cell is
    reported as OOM for the original mask-construction implementation. The
    current full-model parity failure means these rows are not admitted as
    end-to-end or batched throughput results.}
\label{tab:t3-flex}
\IfFileExists{generated/T3_flex.tex}{%
    % Generated by hand from flex timing grid (2026-08-28).
% Source: results/runs.jsonl, logs/timing_flex_grid3.log
\begin{tabular}{rr rr rr r}
\toprule
$N$ & $\ell$ & \multicolumn{2}{c}{FTT (ms)} & \multicolumn{2}{c}{DDI ratio} & System \\
 & & sdpa full & flex ddi\_qf & sdpa & flex & speedup \\
\midrule
3 & 512 & 151 & 147 & 0.98$\times$ & 0.98$\times$ & 1.03$\times$ \\
9 & 512 & 485 & 397 & 1.00$\times$ & 1.05$\times$ & 1.22$\times$ \\
25 & 512 & 1901 & 1066 & 1.01$\times$ & 1.21$\times$ & 1.78$\times$ \\
3 & 2048 & 685 & 553 & 0.98$\times$ & 1.01$\times$ & 1.24$\times$ \\
9 & 2048 & 3247 & 1601 & 1.00$\times$ & 1.25$\times$ & 2.03$\times$ \\
25 & 2048 & 17762 & 4458 & 1.01$\times$ & \textbf{1.81$\times$} & \textbf{3.98$\times$} \\
\bottomrule
\end{tabular}
%
}{%
    \todoreport{T3-flex table}{make timing-flex}%
}
\end{table}

Table~\ref{tab:t3} reports prefill and per-output-token measurements, not
end-to-end latency. End-to-end evidence is limited to the separate corrected
cache harness in Section~\ref{sec:results-caching}.
